\subsection*{Actual subgroup images do not remove the remaining class-two line}

\statement{B.60}{Proved}{A nonzero common class-two image has explicit geometric lifts}
Retain $G=P/\gamma_3P$ and the primitive central vector $w$
from B.59.  The images of the two actual geometric centralizers
$H_T,H_5$ in $G$ have a common subgroup containing
$2\mathbb Z w$.  Their completed images in $\widehat G$
have a common subgroup containing $2\widehat{\mathbb Z}w$.
Consequently restricting the possible line in B.59 to the
individual subgroup images at class two cannot eliminate it.
These are common images in a quotient; no nonzero element
of $\mathcal K$ is constructed.

\begin{proof}
Write $e_1,\ldots,e_{13}$ for the exact meridian basis of
B.56, and let $v_i$ be its ambient degree-one images.
The integral pair-linking matrix from B.58, followed by
the bracket formula (B.59.5), gives
\[
 [v_1,v_5]-[v_2,v_5]+[v_3,v_5]=2w.
                                                               \tag{B.60.1}
\]
The complete $165\times78$ bracket matrix has rank $13$;
adjoining $w$ does not increase its rational rank.
The identity (B.60.1), including the factor two, is verified
in ordinary exact arithmetic by
\nolinkurl{next_B59_image/check_bracket_span.py} and recorded
with the complete matrix in
\nolinkurl{next_B59_image/bracket_span.json}.
It supplies the actual finite word
\[
                  a=[e_1e_2^{-1}e_3,e_5]\in H_T,
              \qquad a\bmod\gamma_3P=2w.                    \tag{B.60.2}
\]

For the cyclic centralizer, the prescribed permutation $D$
fixes labels $1,3$ and has cycles
\[
 (2,6,12,9,4),\qquad(5,7,11,10,8).
\]
Put its two fixed poles at $0,\infty$, and write the two
regular orbits as $\{\zeta^i\}$ and $\{u\zeta^i\}$,
in the displayed cyclic orders, where $\zeta$ is a primitive
fifth root of unity.  Here
$u\in\mathbb P^1\setminus(\{0,\infty\}\cup\mu_5)$.
Based positive meridians $f_0,f_1,\ldots,f_5$ about
$0,1,\zeta,\ldots,\zeta^4$ give actual elements of $H_5$.
Indeed the paths preserve the given cyclic action, and
isotopy extension on its quotient lifts to an equivariant
isotopy fixing the two poles.  The resulting endpoint
mapping classes are pure and commute with $D$.
The choice of paths to the punctures does not affect their
degree-one classes or the following degree-two calculation.

The pair-linking vector at $u=0$ has coefficient one
on every pair of moving orbit points and on each pair
of a moving point with the zero pole.  At $u=\zeta^r$,
it has coefficient one on the five pairs
$(\zeta^{j+r},u\zeta^j)$ and zero elsewhere.
Apply B.58's integral normalization matrix $A$ to these
six vectors, obtaining $q_0,\ldots,q_5\in\mathbb Z^{54}$.
The exact program
\nolinkurl{next_B59_image/check_cyclic_bracket_span.py}
checks that these vectors are fixed by $D$ and independent,
that their fifteen bracket columns are independent, and that
\begin{align*}
 w={}&-2[q_0,q_1]-[q_0,q_2]-[q_0,q_3]
       -[q_0,q_4]-[q_0,q_5]\\
    &+[q_1,q_2]+[q_1,q_5].                                  \tag{B.60.3}
\end{align*}
All entries, labels and coefficients are supplied in
\nolinkurl{next_B59_image/cyclic_bracket_span.json}.
The independent root review recomputes both $165$-coordinate
identities directly from the original degree-one columns,
without using either stored bracket matrix; its passing
record is \nolinkurl{records/root_review_B60.txt}.
Thus the finite word
\[
 b_0=[f_0,f_1]^{-2}[f_0,f_2]^{-1}[f_0,f_3]^{-1}
     [f_0,f_4]^{-1}[f_0,f_5]^{-1}[f_1,f_2][f_1,f_5]
                                                               \tag{B.60.4}
\]
has image $w$ in $G$.  In particular $b=b_0^2\in H_5$
has the same image $2w$ as $a$.

These calculations concern the specified actual embeddings.
An identification of the cyclic family with its geometric
base configuration uses a path of labelled configurations;
changing this path conjugates the resulting mapping classes
by a pure mapping class.  Such conjugation is trivial on
$\gamma_2P/\gamma_3P$, so it does not change (B.60.3).
The same observation accounts for the meridian base paths.
No assertion that the two finite words coincide in $P$
is made.

The profinite assertion uses the full topology from B.59.
The central lattice of $\widehat G$ is
$\widehat{\mathbb Z}^{165}$, and $w$ has a coordinate equal
to $1$.  Therefore the power maps of the images of $a$ and
$b_0$ have images $2\widehat{\mathbb Z}w$ and
$\widehat{\mathbb Z}w$, respectively, and are injective
on $\widehat{\mathbb Z}$.  Equivalently, the finite coordinate
quotients detect every finite congruence of the corresponding
infinite cyclic subgroups.  These closed cyclic images
are contained in the completed images of the actual
centralizers, proving the assertion with the full profinite
topology.

Finally, the particular tetrahedral lift $a$ is not in the
full rigid-seven kernel.  Substitution of its source word
$(1,-2,3,5,-3,2,-1,-5)$ into the certified thirteen
six-strand images gives a nonidentity unreduced Burau
matrix over $\mathbb F_{101}$ at $t=2$: its $(1,1)$ entry
is $27$.  The defining braid relations and inverses are
verified in
\nolinkurl{next_B59_image/check_two_w_rigid.py}; the actual
output and complete matrix are retained.
This source word is a commutator, so its braid exponent
sum is zero.  A central full-twist power would therefore
be the identity, contrary to that finite matrix witness.
By B.32, this explicitly rules out $a\in\mathcal K$.
\end{proof}
